\chapter*{Resumen}

\section*{\tituloPortadaVal}

El presente Trabajo Fin de Grado aborda el diseño y desarrollo de una herramienta basada en visión por computador para la segmentación de imágenes de resonancia magnética (IRM) con el objetivo de medir volúmenes de un órgano y una lesión en un marco de investigación preclínico.

En colaboración con la ICTS BioImagen Complutense, se ha desarrollado un prototipo funcional que realiza segmentación automática de uno de los casos de uso planteados por la ICTS BioImagen Complutense: la segmentación de lesión isquémica en cerebro de ratón. Usando dos modelos, la herramienta realiza el proceso de segmentación sobre cada uno de los cortes, primero sobre el cerebro completo y segundo sobre la zona dañada. Después, con los datos de segmentación de cada región se calculan los volúmenes del cerebro total y del área dañada.

Debido a la escasez de datos anotados, habitual en el ámbito biomédico, la metodología se ha fundamentado en el uso de \textit{transfer learning} sobre una arquitectura U-Net. Asimismo, para facilitar el análisis cualitativo de las predicciones y aumentar la transparencia del prototipo, se han incorporado técnicas de Inteligencia Artificial Explicable (XAI), como Grad-CAM e \textit{Integrated Gradients}.
\section*{Palabras clave}

\noindent Imagen de Resonancia Magnética (IRM), Isquemia cerebral, Segmentación médica, Region of Interest (ROI), Inteligencia Artificial (IA), Redes Neuronales Convolucionales (CNN), U-Net, Inteligencia Artificial Explicable (XAI), Visión por computador.



